% !TeX root = ../papers/nc-p3.tex
\section{階数14の四中心商と六次の道空間}
\label{sec:rank-fourteen-central-cut}
\begin{theorem}\label{thm:rank-fourteen-central-cut}
標数零の階数14の二成分上で、各通常幾何点において、文字 $0,1$ から二本ずつ
四次中心元を選び、全四始点で商の長さ $0$ から $6$ の次元を
$1,4,10,20,31,40,44$ にできる。特に割る前の長さ6の道空間はどの始点でも次元84である。
以下の計算で生じる $t$ 直線上の22既約閉点の上では、全4768幾何点で
標準Koszul左右AS正則性が成立する。
\end{theorem}
\begin{proof}
前定理の中心行列の自由列を左から数え、文字0の第一・第三自由変数と
文字1の二自由変数に対応する四つの核ベクトルを取る。
\path{computations/verify_rank_fourteen_central_cut.py} は始点 $p=0,1$ ごとに
語基底の右に一文字を付ける。長さ $n$、文字 $h$ の列は、末尾 $b$ ごとに
長さ $n-1$、文字 $h+b$ の残存語を並べる（加法は四元群のもの）。
長さ $n-2$ の残存語の後ろに各二次関係を付け、最後から二番目の文字までを
直前の射影で展開して一行を得る。長さ4では選んだ中心式も行に加える。
これで中心式の全挿入を扱える。実際、四次中心元の各位置の成分は直前の矢と
交換でき、途中で始まる四次関係は先頭へ移してから右に矢を付けて得られる。

$(\ell,t,i,Y)=(101,2,10,38)$ の剰余点で、短い行から順に独立な行を選ぶ。
選んだ行について、ピボット列・行列式・自由列への射影は前定理の関数体で改めて計算する。
以前の行の関係を使い、残った非零行の最左列を順に消せば、
ピボットの積は選択小行列式に符号を掛けたものである。
この行列式が可逆なら、Cramerの公式で新しい自由列は全ての語を張る。
各段で一部の独立行だけを使っても、得られる語数は商次元の上界となる。
各始点で長さ $0$ から $6$ の上界は定理の列になる。
関係・中心式の分母、経路基底の極、および選択小行列式のノルムの分子と分母から
非通常因子を除くと、相異なる有理既約因子の次数は
\[
\begin{gathered}
8,8,8,14,16,16,20,22,24,24,34,34,\\
52,52,76,100,106,108,112,114,116,128 .
\end{gathered}
\]
この集合の外では、二位置移動と $Y$ の符号交換により残る二始点でも
同じ上界を与える切断が存在する。各始点でこの上界を満たす選択は
$\operatorname{Gr}(2,Z_{4,0})$ の非空開集合を含む（文字1では空間全体を取る）。
Grassmann多様体は既約なので、四つの開集合の共通部分から一組を取れる。

22因子については、それぞれの単純な剰余根を検算コードに指定した。
各根で $i$ と $s$ の全四組、全四始点について四中心商の長さ13までを再計算し、
命題\ref{prop:fixed-residual-regularity}の列を得る。従って持上げと共役により、
次数の和1192の四倍の幾何点を全次数で認証した。

最後に、割る前の次元には一般正則点から $\binom{n+3}{3}$ の下界がある。
各始点の商の斉次基底を持ち上げれば、それらは次数4の四変数多項式環上で
道空間を生成する。長さ6までは、上の商次元の上界を
$(1-u^4)^{-4}$ に掛けると各係数がこの下界に一致する。
従ってこの範囲では商と道空間の双方で等号が成立する。
\end{proof}
